\documentclass[ngerman,11pt]{article}

\usepackage[utf8]{fontenc}
\usepackage[left=1.5cm, right=1.5cm, top=1cm, bottom=1cm]{geometry}
\usepackage{babel}
\usepackage{listings}
\usepackage{color}
\usepackage{hyperref}
\usepackage{csquotes}
\MakeOuterQuote{"}

\definecolor{dkgreen}{rgb}{0,0.6,0}
\definecolor{gray}{rgb}{0.5,0.5,0.5}
\definecolor{mauve}{rgb}{0.58,0,0.82}


\pagenumbering{gobble}
\renewcommand{\lstlistingname}{Schnipsel}

\lstset{frame=tb,
    language=Java,
    aboveskip=3mm,
    belowskip=3mm,
    showstringspaces=false,
    captionpos=b,
    columns=flexible,
    basicstyle={\small\ttfamily},
    numbers=left,
    numberstyle=\tiny\color{gray},
    keywordstyle=\color{blue},
    commentstyle=\color{dkgreen},
    stringstyle=\color{mauve},
    breaklines=true,
    breakatwhitespace=true,
    tabsize=3,
    inputencoding=utf8
}

% Packages
\usepackage{amsmath}

% Document
\begin{document}
{\huge \textbf{Informatik - Programmieren mit Java - Schleifen}}

    \setlength{\parindent}{0}


    \section{Wiederholung - \lstinline{boolean} und \lstinline{if-else}}

    Eine Variable vom Typ \lstinline{boolean} ist entweder \lstinline{true} oder \lstinline{false} (wahr oder falsch).
    Mithilfe von \lstinline{if-else} können wir steuern, wie sich das Programm in bestimmten
    Fällen verhalten soll.

    Beispiele, wie wir booleans erzeugen können:

    \begin{lstlisting}[caption=\lstinline{Vergleiche}]
int x = 3;
int y = 8;
boolean b1 = x < y; // ist x kleiner als y?
boolean b2 = x > y; // ist x groesser als y?
boolean b3 = x <= y; // ist x kleiner ODER gleich y?
boolean b4 = x >= y; // ist x groesser ODER gleich y?
boolean b5 = x == y; // ist x gleich y?
    \end{lstlisting}

    \begin{lstlisting}[caption=\lstinline{if-else},label={lst:x}]
public class Main {
    public static void main(String[] args) {
        int jahr = 2027;
        boolean istInVergangenheit = jahr < 2024;
        if(istInVergangenheit) {
            System.out.println("Willkommen in der Vergangenheit!");
        } else {
            System.out.println("Dies ist die Zukunft!");
        }
    }
}
    \end{lstlisting}

    \textbf{Aufgabe 1}: Tippe das Programm aus Schnipsel~\ref{lst:x} ab, führe es aus und überprüfe, ob du die Ausgabe nachvollziehen kannst. \\

    \textbf{Aufgabe 2}: Ändere das Programm wie folgt: Wenn die Variable \lstinline{jahr} dem aktuellen Jahr entspricht (2024),
    soll die Ausgabe \lstinline{"Hallo, Gegenwart!"} sein.
    Ansonsten soll die Ausgabe so sein wie vorher
    (es sollen also insgesamt 3 verschiedene Fälle unterschieden werden).

    \section{\lstinline{while}-Schleife}

    Wie würdest du ein Programm schreiben, das dreimal \lstinline{"Maracuja!"} auf der Konsole ausgibt?
    Vielleicht in etwa so:

    \begin{lstlisting}
System.out.println("Maracuja!");
System.out.println("Maracuja!");
System.out.println("Maracuja!");
    \end{lstlisting}

    Okay, sollte klappen.
    Wie sähe es denn aus mit 1000 Mal?

    (Nein, bitte nicht 1000 Mal die selbe Zeile einfügen!) \\

    Um Teile unseres Programms mehrfach auszuführen, können wir \textit{Schleifen} verwenden.
    Hier schauen wir uns zunächst die \lstinline{while}-Schleife an (while ist englisch für \textit{solange}). \\

    \newpage
    Auf deutsch könnten wir formulieren: \textit{"Solange eine Bedingung wahr ist, tue dies"}

    Der Java Code für dieses Verhalten sieht wie folgt aus:

    \begin{lstlisting}[caption=\lstinline{while}]
while(Bedingung) {
    // "dies"
}
    \end{lstlisting}

    Jedes mal, bevor "dies" (also der Code innerhalb des \lstinline{while}-Blocks) ausgeführt wird, wird die Bedingung (erneut) überprüft.
    Ist sie (immer noch) wahr, wird der Code (nochmal) ausgeführt.
    Sonst fährt die Ausführung nach der Schleife fort. \\

    Kommen wir zurück zu den Maracujas:

    \begin{lstlisting}[caption=\lstinline{while}, label={while}]
public class Main {
    public static void main(String[] args) {
        int zaehler = 0;
        while(zaehler < 1000) {
            System.out.println("Maracuja!");

            // Wichtig: Zaehlvariable anpassen, sonst bleibt die Bedingung oben immer wahr, und die Schleife laeuft fuer immer!
            zaehler = zaehler + 1;
        }
    }
}
    \end{lstlisting}

    \textbf{Aufgabe 3}: Tippe das Programm aus Schnipsel \ref{while} ab, führe es aus, und überprüfe, ob du das Verhalten nachvollziehen kannst. \\

    \textbf{Aufgabe 4}: Schreibe ein Programm, das die Zahlen von 1 bis 100 ausgibt. \\

    \textbf{Aufgabe 5}: Schreibe ein Programm, das die Zahlen von 1 bis 100 ausgibt, wobei hinter den Zahlen 1 bis 10 ein \lstinline{"!"} ausgegeben wird.

    \textit{Tipp}: Wir können Strings und Zahlen zu neuen Strings kombinieren: \lstinline{"hallo"  + 5} ergibt den String \lstinline{"hallo 5"}. \\

    \textbf{Aufgabe 5.1}: Falls du in Aufgabe 5 mehr als eine \lstinline{while}-Schleife verwendet hast, finde eine Lösung, welche nur mit einer \lstinline{while}-Schleife auskommt. \\

    \newpage
    \section{Muster "zeichnen"}

    \begin{lstlisting}[caption={mehrzeilige Ausgabe}, label={println}]
public class Main {
    public static void main(String[] args) {
        System.out.println("Tick");
        System.out.println("Trick");
        System.out.println("Track");
    }
}
    \end{lstlisting}
    \begin{lstlisting}[caption={einzeilige Ausgabe}, label={print}]
public class Main {
    public static void main(String[] args) {
        System.out.print("Tick");
        System.out.print("Trick");
        System.out.print("Track");
    }
}
    \end{lstlisting}

    \textbf{Aufgabe 6}: Vergleiche die Ausgaben der Programme aus Schnipsel \ref{println} und Schnipsel \ref{print}.
    Wo ist der Unterschied? \\

    \textbf{Aufgabe 7}: Schreibe ein Programm, welches 10 Mal direkt hintereinander (in einer Zeile, ohne Leerzeichen zwischendurch) \lstinline{"Erdbeere"} ausgibt. \\

    \textbf{Aufgabe 8}: Schreibe ein Programm, welches wie folgt arbeitet:
    \begin{itemize}
        \item Es gibt eine Variable vom Typ \lstinline{int} namens \lstinline{anzahl}
        \item Es soll folgendes ausgegeben werden: \lstinline{1*2*3*4*5*....} (bis die Zahl \lstinline{anzahl} erreicht ist).
    \end{itemize} \\

    \vspace{1cm}

    Lasst uns nun tatsächlich Muster zeichnen! \\
    Das folgende Programm gibt ein Quadrat aus:

    \begin{lstlisting}[caption=\lstinline{Muster: Quadrat}, label={quadrat}]
public class Main {
    public static void main(String[] args) {
        int groesse = 7;

        int zaehlerZeile = 0;
        while (zaehlerZeile < groesse) {
            zaehlerZeile = zaehlerZeile + 1;

            int zaehlerSpalte = 0;
            while (zaehlerSpalte < groesse) {
                zaehlerSpalte = zaehlerSpalte + 1;

                System.out.print("* ");
            }

            System.out.println();
        }
    }
}
    \end{lstlisting}

    \textbf{Aufgabe 9}: Ändere den Code für das Quadrat so ab, dass ein Rechteck ausgegeben wird (Länge und Breite sollen verschieden sein können). \\

    \textbf{Aufgabe 10}: Ändere den Code für das Quadrat so ab, dass stattdessen ein Dreieck ausgegeben wird.

    \textit{Tipp}: Überlege dir, wie du dafür die Bedingung der inneren \lstinline{while}-Schleife anpassen kannst. \\

    \textbf{Bonusaufgabe}: Schreibe ein Programm, welches einen Kreis ausgibt.

    \textit{Tipp}: Satz des Pythagoras

    \newpage

    \section{Musterlösung}

    \textbf{Aufgabe 1} Die Ausgabe sollte "Willkommen in der Vergangenheit" sein, solange das jahr kleiner ist als 2024 (also bis einschließlich 2023). \\
    Für mindestens den Wert 2024 (inklusive der vorgegebenen 2027) sollte die Ausgabe "Dies ist die Zukunft!" sein. \\

    \textbf{Aufgabe 2}

    \begin{lstlisting}
public class Main {
    public static void main(String[] args) {
        int jahr = 2027;
        boolean istInVergangenheit = jahr < 2024;
        if(istInVergangenheit) {
            System.out.println("Willkommen in der Vergangenheit!");
        } else {
            // Hier kann man alternativ erst eine Variable wie vorher definieren, und diese im if verwenden:
            // boolean istGegenwart = jahr == 2024;
            if(jahr == 2024) {
                System.out.println("Hallo, Gegenwart!");
            } else {
                System.out.println("Dies ist die Zukunft!");
            }
        }
    }
}
    \end{lstlisting} \\

    \textbf{Aufgabe 3}: Es sollten 1000 Zeilen mit dem Text "Maracuja!" ausgegeben werden. \\

    \textbf{Aufgabe 4}:

    \begin{lstlisting}
public class Main {
    public static void main(String[] args) {
        int zahl = 1;
        while(zahl <= 100) {
            System.out.println(zahl);
            zahl = zahl + 1;
        }
    }
}
    \end{lstlisting} \\

    \textbf{Aufgabe 5 bzw 5.1}:

    \begin{lstlisting}
public class Main {
    public static void main(String[] args) {
        int zahl = 1;
        while(zahl <= 100) {
            if(zahl <= 10) {
                System.out.println(zahl + "!");
            } else {
                System.out.println(zahl);
            }
            zahl = zahl + 1;
        }
    }
}
    \end{lstlisting} \\

    \textbf{Aufgabe 6}: Mit \lstinline{System.out.println} wird der Ausgabe ein Zeilenumbruch angefügt, nachfolgende Ausgaben
    geschehen also in einer neuen Konsolenzeile.
    Mit \lstinline{System.out.print} passiert das nicht, die Ausgaben werden alle in einer Zeile hintereinander gereiht.
    \textit{Bemerkung}: Man kann den Schülern hier den Hinweis geben, dass der Befehl \lstinline{System.out.println()} nur einen Zeilenumbruch produziert. \\

    \textbf{Aufgabe 7}

    \begin{lstlisting}
public class Main {
    public static void main(String[] args) {
        int zahl = 1;
        while(zahl <= 10) {
            System.out.print("Erdbeere");
            zahl = zahl + 1;
        }
    }
}
    \end{lstlisting} \\

    \textbf{Aufgabe 8}:

    \begin{lstlisting}
public class Main {
    public static void main(String[] args) {
        int anzahl = 50;
        int zaehler = 0;
        while(zaehler <= anzahl) {
            System.out.print(zahl + "*");
        }
    }
}
    \end{lstlisting} \\

    \textbf{Aufgabe 9}:

    \begin{lstlisting}[caption=\lstinline{Muster: Rechteck}]
public class Main {
    public static void main(String[] args) {
        int hoehe = 6;
        int breite = 10;

        int zaehlerZeile = 0;
        while (zaehlerZeile < hoehe) {
            zaehlerZeile = zaehlerZeile + 1;

            int zaehlerSpalte = 0;
            while (zaehlerSpalte < breite) {
                zaehlerSpalte = zaehlerSpalte + 1;

                System.out.print("* ");
            }

            System.out.println();
        }
    }
}
    \end{lstlisting}

    \textbf{Aufgabe 10}:

    \begin{lstlisting}[caption=\lstinline{Muster: Dreieck}]
public class Main {
    public static void main(String[] args) {
        int groesse = 7;

        int zaehlerZeile = 0;
        while (zaehlerZeile < groesse) {
            zaehlerZeile = zaehlerZeile + 1;

            int zaehlerSpalte = 0;
            while (zaehlerSpalte < zaehlerZeile) {
                zaehlerSpalte = zaehlerSpalte + 1;

                System.out.print("* ");
            }

            System.out.println();
        }
    }
}
    \end{lstlisting} \\

    \textbf{Bonusaufgabe}

    \begin{lstlisting}
public class Main {
    public static void main(String[] args) {
        int radius = 20;
        int x = 0;
        while(x < radius * 2 + 1) {
            int y = 0;
            while(y < radius * 2 + 1) {
                int xRelativZurMitte = x - radius;
                int yRelativZurMitte = y - radius;
                if(xRelativZurMitte * xRelativZurMitte + yRelativZurMitte * yRelativZurMitte <= radius * radius) {
                    System.out.print("O ");
                } else {
                    System.out.print("- ");
                }
                y = y + 1;
            }
            System.out.println();
            x = x + 1;
        }
    }
}
        \end{lstlisting}

    \end{document}
